\documentclass[11pt]{article}
\usepackage[margin=1in]{geometry}
\usepackage{amsmath,amssymb,amsthm}
\usepackage{graphicx}
\usepackage{hyperref}

\newtheorem{theorem}{Theorem}
\newtheorem{corollary}[theorem]{Corollary}
\newtheorem{lemma}[theorem]{Lemma}
\newcommand{\BMO}{\mathrm{BMO}}
\newcommand{\R}{\mathbb R}
\newcommand{\cI}{\mathcal I}

\title{Carleson-Packed Product-BMO Lower Bounds\\for Bi-Commutators}
\author{Run \texttt{fa\_banach\_001}, lane 8}
\date{August 11, 2026}

\begin{document}
\maketitle

\section*{Status and source question}

\textbf{Status:} candidate substantial partial result, likely valid; human
review requested.

Hytönen records in Section 1.6(2), printed page 7, that the diagonal
$L^p\to L^p$ characterization of the bi-commutator
\[
  [T_1,[M_b,T_2]]
\]
remains open after a gap was found in the Ferguson--Lacey theorem and its
extensions \cite{Hy}.  The exact source passage is reproduced below.

\begin{center}
\includegraphics[width=0.96\textwidth]{figures/open_problem_crop.png}
\end{center}

For the Hilbert transforms on $\R^2$, the missing direction is the lower
bound
\[
 \|b\|_{\BMO_{\rm prod}}\lesssim
 \|[H_1,[M_b,H_2]]\|_{L^2\to L^2}.
\]
Airta, Hytönen, Li, Martikainen, and Oikari explain the false
Fourier-quadrant step and prove the surviving lower bound by rectangular BMO
for general nondegenerate kernels \cite{AHLMO}.  This packet upgrades that
bound to product BMO for a new support-packing class, and separately settles
all pure tensor symbols.  The unrestricted problem remains open.

\section*{Idea of the proof}

Rectangular BMO says that after a first-coordinate interval $I$ is fixed, the
remaining coefficients indexed by $J$ form a one-parameter Carleson sequence
with constant at most $B^2|I|$.  For an arbitrary open set $\Omega$, let
$\Omega_I$ collect the $J$ for which $I\times J\subset\Omega$.  The slice
coefficient mass is at most $B^2|I||\Omega_I|$.

The only issue is summing this over $I$.  If the active $I$'s form a Carleson
family, layer cake bounds the sum by the packing constant times $|\Omega|$.
This produces exactly the open-set Carleson inequality defining product BMO.
It identifies arbitrary depth of incomparable rectangles as the obstruction
left in the full problem.

For $b=u\otimes v$, the same layer-cake argument is weighted by the Haar
coefficients of $u$, themselves a Carleson sequence.  Simultaneously, the
bi-commutator factors into the tensor product of the two one-parameter
commutators, giving a dimension-free result without any scale restriction.

\section*{Definitions}

Fix a tensor-product dyadic or Meyer wavelet basis
$\{w_I\otimes w_J:I,J\in\mathcal D\}$, and put
\[
 b_{I,J}=\langle b,w_I\otimes w_J\rangle.
\]
The coefficient forms of dyadic product and rectangular BMO are
\begin{align*}
 \|b\|_{\BMO_{\rm prod}}^2
 &=\sup_{0<|\Omega|<\infty}\frac1{|\Omega|}
   \sum_{I\times J\subset\Omega}|b_{I,J}|^2,\\
 \|b\|_{\BMO_{\rm rect}}^2
 &=\sup_{I_0,J_0}\frac1{|I_0||J_0|}
   \sum_{I\subset I_0,\,J\subset J_0}|b_{I,J}|^2.
\end{align*}
For a standard Meyer system these norms are equivalent to the corresponding
continuous Chang--Fefferman and rectangular norms, with constants depending
only on the fixed system \cite{CF}.

Define the first-coordinate active family
\[
 \cI_b=\{I\in\mathcal D:b_{I,J}\ne0\text{ for some }J\}.
\]
Its packing constant is
\[
 \Lambda_1(\cI_b)=
 \sup_{0<|E|<\infty}\frac1{|E|}
 \sum_{I\in\cI_b,\,I\subset E}|I|,
\]
where $E\subset\R$ ranges over open sets.  Define $\Lambda_2$ symmetrically.

\section*{Packing theorem}

\begin{theorem}[Support-packing upgrade]\label{thm:packing}
If $\Lambda_1(\cI_b)<\infty$, then
\[
 \|b\|_{\BMO_{\rm prod}}
 \le \sqrt{\Lambda_1(\cI_b)}\,
       \|b\|_{\BMO_{\rm rect}}.
\]
The symmetric estimate holds with $\Lambda_2$.  Hence the better of the two
finite packing constants may be used.
\end{theorem}

\begin{proof}
Write $B=\|b\|_{\BMO_{\rm rect}}$.  For fixed $I$ and every dyadic $J_0$,
\[
 \sum_{J\subset J_0}|b_{I,J}|^2
 \le \sum_{I'\subset I,\,J\subset J_0}|b_{I',J}|^2
 \le B^2|I||J_0|.                                      \tag{1}
\]
Thus $\{|b_{I,J}|^2\}_J$ is a one-parameter Carleson sequence with constant
at most $B^2|I|$.

Fix an open $\Omega\subset\R^2$ of finite measure, and let $\Omega_I$ be the
union of all dyadic intervals $J$ such that $I\times J\subset\Omega$.
Decompose $\Omega_I$ into maximal dyadic intervals.  Applying (1) on these
disjoint intervals gives
\[
 \sum_{J:I\times J\subset\Omega}|b_{I,J}|^2
 \le B^2|I||\Omega_I|.                                  \tag{2}
\]

Put $g_I=|\Omega_I|$ and
$E_t=\bigcup\{I\in\cI_b:g_I>t\}$.  Layer cake and the packing hypothesis
give
\begin{align*}
 \sum_{I\in\cI_b}|I|g_I
 &=\int_0^\infty\sum_{I:g_I>t}|I|\,dt\\
 &\le\Lambda_1\int_0^\infty|E_t|\,dt.                  \tag{3}
\end{align*}
If $x\in E_t$, then $x\in I$ for an interval with $|\Omega_I|>t$, and
$\Omega_I$ is contained, up to endpoints, in the vertical section
$\Omega_x$.  Therefore $E_t\subset\{x:|\Omega_x|>t\}$, and Fubini gives
\[
 \int_0^\infty|E_t|\,dt
 \le\int_{\R}|\Omega_x|\,dx=|\Omega|.                  \tag{4}
\]
Summing (2), then using (3)--(4), proves
\[
 \sum_{I\times J\subset\Omega}|b_{I,J}|^2
 \le B^2\Lambda_1|\Omega|.
\]
Taking the supremum over $\Omega$ completes the proof.  The other coordinate
is identical.
\end{proof}

\begin{corollary}[Finite one-coordinate scale complexity]
If the active wavelet coefficients use at most $M$ distinct scales in either
coordinate, then
\[
 \|b\|_{\BMO_{\rm prod}}
 \le\sqrt M\,\|b\|_{\BMO_{\rm rect}}.                  \tag{5}
\]
Consequently, for symmetrically nondegenerate Calderón--Zygmund operators in
the two coordinates,
\[
 \|b\|_{\BMO_{\rm prod}}
 \lesssim\sqrt M\,
 \|[T_1,[M_b,T_2]]\|_{L^2\to L^2}.                    \tag{6}
\]
Together with the standard product-BMO upper bound, (6) is a full
boundedness characterization on this finite-complexity class.
\end{corollary}

\begin{proof}
At each fixed scale, the dyadic intervals contained in an open set are
disjoint, so their lengths sum to at most the measure of the set.  Hence an
$M$-scale active family has packing constant at most $M$.  Apply
Theorem~\ref{thm:packing}; then use the rectangular lower bound of
\cite[Proposition 5.3]{AHLMO}, whose off-support constant is dominated by the
operator norm.
\end{proof}

\section*{Pure tensor theorem}

\begin{theorem}[Rank-one symbols]\label{thm:tensor}
For one-parameter functions $u,v$,
\[
 \|u\otimes v\|_{\BMO_{\rm prod}}
 \asymp \|u\|_{\BMO}\|v\|_{\BMO}.                    \tag{7}
\]
Moreover,
\[
 [T_1,[M_{u\otimes v},T_2]]
   =[T_1,M_u]\otimes[M_v,T_2].                         \tag{8}
\]
In particular, for $T_1,T_2$ having the usual one-parameter commutator
characterization of BMO (including the Hilbert transforms), diagonal
bi-commutator boundedness is equivalent to product BMO on all pure tensor
symbols, with no scale restriction.
\end{theorem}

\begin{proof}
Let $a_I=\langle u,w_I\rangle$ and $c_J=\langle v,w_J\rangle$.
Testing the product-BMO norm on $I_0\times J_0$ and approximating the two
one-parameter suprema proves the lower bound in (7).

For the upper bound, fix $\Omega$ and define $\Omega_I$ as above.  The
Carleson property of $\{|c_J|^2\}$ gives
\[
 \sum_{J:I\times J\subset\Omega}|a_Ic_J|^2
 \le |a_I|^2\|v\|_{\BMO}^2|\Omega_I|.                 \tag{9}
\]
Using layer cake again, now with the Carleson weights $|a_I|^2$, yields
\begin{align*}
 \sum_I|a_I|^2|\Omega_I|
 &=\int_0^\infty\sum_{I:|\Omega_I|>t}|a_I|^2\,dt\\
 &\le\|u\|_{\BMO}^2
   \int_0^\infty|\{x:|\Omega_x|>t\}|\,dt\\
 &=\|u\|_{\BMO}^2|\Omega|.                           \tag{10}
\end{align*}
Equations (9)--(10) prove the upper bound in (7).

Identity (8) follows directly because $T_1$ and $T_2$ act in separate
coordinates.  Operator norms multiply on Hilbert tensor products.  The
one-parameter commutator theorem \cite{CRW} and (7) finish the proof.
\end{proof}

\section*{Verification, upgrade attempts, and limitations}

The proof is analytic and uses no computation.  Eight focused upgrade routes
were checked before promotion: current-literature status, summing all Fourier
quadrants, modulation/reflection, replacing bounded overlap by Carleson
packing, exploiting small rectangular BMO in the old bootstrap, Rademacher
randomization, smooth symbol-frequency projections, and finite tensor-rank
factorization.  The first four produced Theorem~\ref{thm:packing}; the last
route produced Theorem~\ref{thm:tensor}.  The remaining routes fail to control
arbitrarily deep incomparable rectangle families uniformly.

Holmes--Treil--Volberg prove a different dyadic partial theorem under an
alternating-scale condition on the Haar support \cite{HTV}.  That result does
not contain the Carleson-packed support class and does not repair the
continuous Hilbert-transform theorem.  Conversely, this packet does not
contain their infinite alternating-scale class.  The unrestricted diagonal
problem remains open.

A bounded search through August 11, 2026 covered the run's four lightweight
indexes, the local sources of \cite{Hy,HTV}, the current version of
\cite{AHLMO}, and official-arXiv queries using the exact problem and the terms
``product BMO'', ``rectangular BMO'', ``finite scales'', and ``Carleson
family''.  No matching support-packing statement was found.  This is not a
priority claim.

Human review should focus on the equivalence between the fixed wavelet
rectangular norm and the oscillation norm used in \cite{AHLMO}, the maximal
dyadic decomposition of $\Omega_I$, and the weighted layer-cake step in
Theorem~\ref{thm:tensor}.

\small
\begin{thebibliography}{9}
\bibitem{Hy}
T. Hytönen,
\emph{The $L^p$-to-$L^q$ boundedness of commutators with applications to the
Jacobian operator}, arXiv:1804.11167, especially Section 1.6(2).

\bibitem{AHLMO}
E. Airta, T. Hytönen, K. Li, H. Martikainen, and T. Oikari,
\emph{Off-diagonal estimates for bi-commutators},
Int. Math. Res. Not. IMRN 2022, no. 23, 18766--18832;
arXiv:2005.03548.

\bibitem{HTV}
I. Holmes, S. Treil, and A. Volberg,
\emph{Dyadic bi-parameter repeated commutator and dyadic product BMO},
arXiv:2101.00763.

\bibitem{CF}
S.-Y. A. Chang and R. Fefferman,
\emph{A continuous version of duality of $H^1$ with BMO on the bidisc},
Ann. of Math. (2) 112 (1980), 179--201.

\bibitem{CRW}
R. Coifman, R. Rochberg, and G. Weiss,
\emph{Factorization theorems for Hardy spaces in several variables},
Ann. of Math. (2) 103 (1976), 611--635.
\end{thebibliography}

\end{document}
